% Options for packages loaded elsewhere
\PassOptionsToPackage{unicode}{hyperref}
\PassOptionsToPackage{hyphens}{url}
%
\documentclass[
]{article}
\usepackage{amsmath,amssymb}
\usepackage{iftex}
\ifPDFTeX
  \usepackage[T1]{fontenc}
  \usepackage[utf8]{inputenc}
  \usepackage{textcomp} % provide euro and other symbols
\else % if luatex or xetex
  \usepackage{unicode-math} % this also loads fontspec
  \defaultfontfeatures{Scale=MatchLowercase}
  \defaultfontfeatures[\rmfamily]{Ligatures=TeX,Scale=1}
\fi
\usepackage{lmodern}
\ifPDFTeX\else
  % xetex/luatex font selection
\fi
% Use upquote if available, for straight quotes in verbatim environments
\IfFileExists{upquote.sty}{\usepackage{upquote}}{}
\IfFileExists{microtype.sty}{% use microtype if available
  \usepackage[]{microtype}
  \UseMicrotypeSet[protrusion]{basicmath} % disable protrusion for tt fonts
}{}
\makeatletter
\@ifundefined{KOMAClassName}{% if non-KOMA class
  \IfFileExists{parskip.sty}{%
    \usepackage{parskip}
  }{% else
    \setlength{\parindent}{0pt}
    \setlength{\parskip}{6pt plus 2pt minus 1pt}}
}{% if KOMA class
  \KOMAoptions{parskip=half}}
\makeatother
\usepackage{xcolor}
\usepackage[margin=1in]{geometry}
\usepackage{graphicx}
\makeatletter
\newsavebox\pandoc@box
\newcommand*\pandocbounded[1]{% scales image to fit in text height/width
  \sbox\pandoc@box{#1}%
  \Gscale@div\@tempa{\textheight}{\dimexpr\ht\pandoc@box+\dp\pandoc@box\relax}%
  \Gscale@div\@tempb{\linewidth}{\wd\pandoc@box}%
  \ifdim\@tempb\p@<\@tempa\p@\let\@tempa\@tempb\fi% select the smaller of both
  \ifdim\@tempa\p@<\p@\scalebox{\@tempa}{\usebox\pandoc@box}%
  \else\usebox{\pandoc@box}%
  \fi%
}
% Set default figure placement to htbp
\def\fps@figure{htbp}
\makeatother
\setlength{\emergencystretch}{3em} % prevent overfull lines
\providecommand{\tightlist}{%
  \setlength{\itemsep}{0pt}\setlength{\parskip}{0pt}}
\setcounter{secnumdepth}{-\maxdimen} % remove section numbering
\usepackage{bookmark}
\IfFileExists{xurl.sty}{\usepackage{xurl}}{} % add URL line breaks if available
\urlstyle{same}
\hypersetup{
  pdftitle={Simple Model Review},
  pdfauthor={Dina Dinh},
  hidelinks,
  pdfcreator={LaTeX via pandoc}}

\title{Simple Model Review}
\author{Dina Dinh}
\date{August 20, 2026}

\begin{document}
\maketitle

\section{Comparing Model 2 and Model
3}\label{comparing-model-2-and-model-3}

\subsection{Where Model 2 and Model 3 Agree in
Contrasts}\label{where-model-2-and-model-3-agree-in-contrasts}

In this scenario, we see that there's a big difference at final between
the two groups, BUT the within-group change is almost zero for both
groups.

\subsection{Where Model 2 and Model 3 Disgree in
Contrasts}\label{where-model-2-and-model-3-disgree-in-contrasts}

In this scenario, we see that at final, there's not much change between
the two groups BUT the within-group change for Group B is large.

Assumptions:

\begin{enumerate}
\def\labelenumi{\arabic{enumi}.}
\tightlist
\item
  Independent subjects across groups
\end{enumerate}

\begin{itemize}
\tightlist
\item
  One person in group A should not be correlated with a person in group
  B.
\end{itemize}

\begin{enumerate}
\def\labelenumi{\arabic{enumi}.}
\setcounter{enumi}{1}
\tightlist
\item
  Change scores are approximately normally distributed within each group
\end{enumerate}

\begin{itemize}
\tightlist
\item
  This does not mean the original outcome at time1 or time2 must each be
  normal. It is the change score that matters most.
\end{itemize}

\begin{enumerate}
\def\labelenumi{\arabic{enumi}.}
\setcounter{enumi}{2}
\tightlist
\item
  Equal variances only if you use the pooled t-test
\end{enumerate}

\begin{itemize}
\tightlist
\item
  If you use the default Welch t-test in R, equal variances are not
  required.
\end{itemize}

\begin{enumerate}
\def\labelenumi{\arabic{enumi}.}
\setcounter{enumi}{3}
\tightlist
\item
  Outcome is continuous
\end{enumerate}

\begin{itemize}
\tightlist
\item
  The change score should be roughly quantitative/continuous.
\end{itemize}

\begin{enumerate}
\def\labelenumi{\arabic{enumi}.}
\setcounter{enumi}{4}
\tightlist
\item
  No major outliers in the change scores
\end{enumerate}

\begin{itemize}
\tightlist
\item
  Extreme outliers can distort the mean and the t-test.
\end{itemize}

Assumptions:

\begin{enumerate}
\def\labelenumi{\arabic{enumi}.}
\tightlist
\item
  Correct mean structure
\end{enumerate}

\begin{itemize}
\tightlist
\item
  The model for the mean should be appropriate:

  \begin{itemize}
  \tightlist
  \item
    group effect
  \item
    time effect
  \item
    group × time interaction
  \end{itemize}
\end{itemize}

\begin{enumerate}
\def\labelenumi{\arabic{enumi}.}
\setcounter{enumi}{1}
\tightlist
\item
  Independence between subjects
\end{enumerate}

\begin{itemize}
\tightlist
\item
  Different subjects should be independent.
\end{itemize}

\begin{enumerate}
\def\labelenumi{\arabic{enumi}.}
\setcounter{enumi}{2}
\tightlist
\item
  Correlation within subjects is handled correctly enough
\end{enumerate}

\begin{itemize}
\tightlist
\item
  The mixed model accounts for repeated observations within subject
  through the random effect or covariance structure.
\end{itemize}

\begin{enumerate}
\def\labelenumi{\arabic{enumi}.}
\setcounter{enumi}{3}
\tightlist
\item
  Residuals are approximately normal
\end{enumerate}

\begin{itemize}
\tightlist
\item
  More specifically, the residuals and random effects are often assumed
  approximately normal for standard inference in linear mixed models.
\end{itemize}

\begin{enumerate}
\def\labelenumi{\arabic{enumi}.}
\setcounter{enumi}{4}
\tightlist
\item
  Homoscedasticity of residuals
\end{enumerate}

\begin{itemize}
\tightlist
\item
  Residual variance should be reasonably similar across fitted
  values/groups/times unless you model different variances.
\end{itemize}

\begin{enumerate}
\def\labelenumi{\arabic{enumi}.}
\setcounter{enumi}{5}
\tightlist
\item
  Outcome is continuous
\end{enumerate}

\begin{itemize}
\tightlist
\item
  For standard linear mixed models.
\end{itemize}

\begin{enumerate}
\def\labelenumi{\arabic{enumi}.}
\setcounter{enumi}{6}
\tightlist
\item
  No extreme influential outliers
\end{enumerate}

\section{ANCOVA}\label{ancova}

At the final timepoint, do treatment and control differ after adjusting
for baseline difference? -\textgreater{} model2 answers this

What is the relationship among the models 2-4? What do they mean and
what are they assuming? For models 2-4, what is the null hypothesis for
the group assuming? Which models 1-4 should we use and in what cases?

model1 : answers if there the a difference between the two groups at
final - null hypothesis: the two group's means are equal to each other
at final - use when only care about final difference between two groups

model2: answers if the two groups differ after adjusting for baseline
difference - null hypothesis: the two group's means are equal to each
other at final after adjusting for baseline - use when want to compare
raw mean changes

model3: answers if the raw change from baseline to final differs between
the two groups - use when want to present changes scores

model4: is the same as model2 except the coefficeient for Baseline is 1
less in model4 - change = \(\beta_0\) + \(\beta_1\)Baseline +
\(\beta_2\)Group - Final - Baseline = \(\beta_0\) + \(\beta_1\)Baseline
+ \(\beta_2\)Group - Final = \(\beta_0\) + \(\beta_1\)Baseline +
Baseline + \(\beta_2\)Group - Final = \(\beta_0\) + (\(\beta_1\) +
1)*Baseline + \(\beta_2\)Group - We can see the only difference in
model2 and model4 mathematically is the +1 - use when baseline
differences between the two groups and want to have outcome variable as
changes scores

\end{document}
